\subsubsection{原则：原理在卷二，本章是 检查清单}

曝光三角、卷帘快门、噪声等原理见卷二「相机与成像」。本章只列赛场上最常踩坑的实践要点。

\begin{example}[建立一次检查记录]
准备一张表，写入日期、相机序列号、镜头、分辨率、曝光、增益、触发方式与标定文件。没有设备时填写“未核对”并注明需查的配置来源。检查：这张表应让另一位同学知道哪些条件已确认，而不是只留下“相机正常”。
\end{example}

\begin{exercise}
\label{v3-ex-camera-1}
检查表中只写“海康相机，默认参数”，为什么不能复现实验？列出四项缺失信息。
\end{exercise}

\subsubsection{卷帘 vs 全局}

\begin{itemize}
    \item \textbf{卷帘快门（Rolling Shutter）}：高速相对运动时易倾斜、撕裂，影响角点与检测稳定性；
    \item \textbf{全局快门（Global Shutter）}：运动场景更稳，成本与选型受限。
\end{itemize}
自瞄对抗、云台快速扫场时，快门类型与曝光时间必须一起评估：曝光过长同样拖影。

\begin{example}[分开观察畸变与拖影]
在已有的快速扫场图像中，分别检查竖直边缘是否整体倾斜、单个灯条边缘是否沿运动方向拉长。前者提示逐行曝光时序效应，后者提示曝光期间运动模糊。记录现象，再用固定其他条件、缩短曝光的对照图检查拖影是否减轻。没有对照图时，只记录候选原因。
\end{example}

\begin{exercise}
\label{v3-ex-camera-2}
全局快门图像仍有明显拖影，能否据此判定相机实际是卷帘快门？
\end{exercise}

\subsubsection{分辨率、帧率与「商家高清」}

\begin{itemize}
    \item 宣传「高清」不等于赛场可用分辨率；
    \item 帧率常与分辨率绑定：高帧可能对应较低分辨率；
    \item 端到端延迟应覆盖约定采样时刻到执行生效之间的读出、传输、排队、处理和执行等待；各阶段可能重叠，应以逐帧时间戳实测，不能用帧率的倒数代替。提高分辨率可能增加处理负担。
\end{itemize}
选型时写清：目标分辨率、最低可接受帧率、接口（USB3/GigE）、黑白/彩色、触发方式。

\begin{example}[用时间戳区分两个指标]
教学样本的连续采样时刻为 0、10、20 ms，对应结果输出时刻为 35、45、55 ms。采样间隔 10 ms，对应 100 帧每秒；逐帧延迟为 35 ms。检查：流水线可每 10 ms 出一帧结果，但每个结果仍反映 35 ms 前的观测。
\end{example}

\begin{exercise}
\label{v3-ex-camera-3}
采样时刻为 0、20、40 ms，结果输出为 50、70、90 ms，采样帧率和逐帧延迟分别是多少？
\end{exercise}

\subsubsection{曝光 / 增益 / 触发 检查清单}

上机前逐项确认：
\begin{enumerate}
    \item 灯条不过曝成一片、暗部仍可见结构；
    \item 增益过高引入的噪声是否拖垮阈值/网络；
    \item 是否已验证曝光策略；自动曝光若拉长曝光或改变灰度分布，可能影响高速运动检测，可锁定曝光或限制自动范围并作对照；
    \item 硬件触发 / 同步是否与电控时间戳约定一致；
    \item 镜头焦距与工作距离是否匹配装甲成像尺寸；
    \item 内参与畸变是否在本机标定，配置是否写对序列号。
\end{enumerate}

雷达俯视与车载自瞄的工作距离、视场（FOV）不同，\textbf{禁止}共用一套未经验证的曝光参数。

\begin{example}[一次只改曝光]
用同一段场景准备两组图像，固定镜头、增益和分辨率，仅改变曝光；记录灯条是否饱和、边缘拖影和暗部可见性。选择兼顾检测与运动条件的一组，保存配置和样图。检查：没有统一的“越亮越好”；缺少设备时可对已有对照图作同样记录。
\end{example}

\begin{note}
切换曝光、触发或采集模式前先保存配置；上机检查时保持执行器处于约定的安全状态。快门类型原理可查厂家说明：\url{https://docs.baslerweb.com/electronic-shutter-types}；实际参数名称与范围以本机型号手册为准。
\end{note}

\begin{conclusion}
\begin{itemize}
  \item 快门类型决定各行曝光的时序关系。
  \item 曝光时长影响拖影和亮度。
  \item 帧率表示采样吞吐，端到端延迟表示某次观测到执行的耗时。
\end{itemize}
\end{conclusion}

\begin{exercise}
\label{v3-ex-camera-4}
换了焦距或裁剪区域后，除了曝光还应核对哪项几何配置？为什么？
\end{exercise}

\subsubsection{工业相机与 SDK*}

赛场常用厂家软件开发工具包（Software Development Kit，SDK，如海康相机 SDK）按序列号打开，而不是笼统 \texttt{VideoCapture(0)}。
驱动、权限、USB 带宽不足时会丢帧——先查系统层再查算法。

\begin{example}[确认打开的是哪台相机]
在厂家查看器中只启动采集预览，读取序列号并与配置文件对照；保存一帧图像，记录分辨率与时间信息，再退出查看器释放设备供程序使用。检查：预览可用仅证明这一采集路径可用，正式程序仍需验证权限、像素格式和时间戳。
\end{example}

\begin{exercise}
\label{v3-ex-camera-5}
查看器能出图而算法程序打不开设备，先检查什么？
\end{exercise}
